\section{Periodicity meets Pell's equation}
\label{sec:pell}

Fix a nonsquare integer $D\ge2$. \emph{Pell's equation} asks for positive
integers $p,q$ such that
\[
 p^2-Dq^2=1.
\]
We will also study the equation with right-hand side $-1$. There are two
questions to settle: why must a solution appear among the convergents of
$\sqrt D$, and which convergents actually work? An approximation estimate
answers the first; an exact identity answers the second.

\subsection*{Every Pell solution is a convergent}
Suppose $p^2-Dq^2=\pm1$, with $p,q>0$. Any common divisor of $p,q$ would
divide $1$, so $p/q$ is reduced. Factoring a difference of squares gives
\begin{equation}\label{eq:pell-error}
 \abs{\sqrt D-\frac pq}
 =\frac{1}{q^2(\sqrt D+p/q)}
 <\frac{1}{2q^2}.
\end{equation}
For the strict inequality, note that
$(p/q)^2\ge D-1\ge1$, while $\sqrt D>1$.
Legendre's criterion now makes $p/q$ a convergent of $\sqrt D$.
This does not yet prove that a solution exists: we still need a convergent
whose difference of squares is exactly $1$ or $-1$.

\subsection*{The palindrome in a square-root period}
Write $a_0=\floor{\sqrt D}$. The first complete quotient is
\[
 x_1=\frac1{\sqrt D-a_0}>1,
 \qquad x_1'=-\frac1{\sqrt D+a_0}\in(-1,0).
\]
Galois's theorem therefore makes $x_1$ purely periodic. Let $L$ be its
\emph{least} period length, so that
$\sqrt D=[a_0;\overline{a_1,\ldots,a_L}]$.

\begin{theorem}[The square-root palindrome]
\label{thm:sqrt-palindrome}
The last digit of the least period is $a_L=2a_0$, and
\[
 a_j=a_{L-j}\qquad(1\le j\le L-1).
\]
Thus
\[
 \boxed{\quad
 \sqrt D=[a_0;\overline{a_1,\ldots,a_{L-1},2a_0}],
 \quad (a_1,\ldots,a_{L-1})\text{ is a palindrome}.
 \quad}
\]
For $L=1$ the palindrome is the empty block.
\end{theorem}
\begin{proof}
Apply Corollary~\ref{cor:period-reversal} to $x_1$:
\[
 -\frac1{x_1'}=\sqrt D+a_0
   =[\overline{a_L,a_{L-1},\ldots,a_1}].
\]
Its integer part is $2a_0$, so $a_L=2a_0$. On the other hand,
\[
 \sqrt D+a_0=2a_0+\frac1{x_1}
   =[2a_0;\overline{a_1,\ldots,a_L}].
\]
These are two expansions of the same irrational number. The
continued-fraction algorithm determines its digits uniquely. Comparing
the next $L-1$ digits gives
$(a_1,\ldots,a_{L-1})=(a_{L-1},\ldots,a_1)$.
\end{proof}

For example,
\[
 \sqrt{23}=[4;\overline{1,3,1,8}],\qquad
 \sqrt{14}=[3;\overline{1,2,1,6}],\qquad
 \sqrt2=[1;\overline2].
\]
The final $2a_0$ marks the end of the symmetric block; it is not itself
part of the palindrome asserted by the theorem.

There is a useful consequence for the integer states
$x_n=(P_n+\sqrt D)/Q_n$ of
\eqref{eq:quadratic-state-recurrence}, starting with $(P_0,Q_0)=(0,1)$.
Since $x_{L+1}=x_1$,
\[
 x_L=a_L+\frac1{x_{L+1}}=2a_0+(\sqrt D-a_0)=\sqrt D+a_0.
\]
Consequently $(P_L,Q_L)=(a_0,1)$, and the same state occurs at every
positive multiple of $L$. Conversely, if $n\ge1$ and $Q_n=1$, reducedness
gives $-1<P_n-\sqrt D<0$, forcing $P_n=a_0$. Thus $x_n=x_L$.
The $L$ states in a least cycle are distinct, so
\begin{equation}\label{eq:pell-state-return}
 \boxed{\qquad Q_n=1\quad\Longleftrightarrow\quad L\divides n
 \qquad(n\ge1).\qquad}
\end{equation}

\subsection*{An exact identity for every convergent}
\begin{lemma}[The complete quotient measures the Pell remainder]
\label{lem:pell-remainder}
For every $n\ge1$,
\begin{equation}\label{eq:pell-remainder}
 \boxed{\quad p_{n-1}^2-Dq_{n-1}^2=(-1)^nQ_n.\quad}
\end{equation}
\end{lemma}
\begin{proof}
Keep the entire tail in the matrix product:
\[
 \sqrt D
 =\frac{p_{n-1}x_n+p_{n-2}}{q_{n-1}x_n+q_{n-2}},
 \qquad x_n=\frac{P_n+\sqrt D}{Q_n}.
\]
Substitute the second expression into the first and multiply through.
Because $\sqrt D$ is irrational, the rational coefficients and the
coefficients of $\sqrt D$ must agree separately. They give
\begin{align*}
 p_{n-1}&=P_nq_{n-1}+Q_nq_{n-2},\\
 Dq_{n-1}&=P_np_{n-1}+Q_np_{n-2}.
\end{align*}
Multiply the first equality by $p_{n-1}$ and the second by $q_{n-1}$,
then subtract:
\begin{align*}
 p_{n-1}^2-Dq_{n-1}^2
 &=Q_n(p_{n-1}q_{n-2}-p_{n-2}q_{n-1})\\
 &=(-1)^nQ_n.
\end{align*}
The last step is the determinant of the product of the $n$ digit
matrices $D(a_0),\ldots,D(a_{n-1})$. The initial conventions
$p_{-1}=1,q_{-1}=0$ also cover $n=1$.
\end{proof}

\begin{theorem}[Exactly which convergents solve Pell]
\label{thm:pell-convergents}
For each $m\ge1$,
\begin{equation}\label{eq:pell-period}
 p_{mL-1}^2-Dq_{mL-1}^2=(-1)^{mL}.
\end{equation}
These are all positive integer solutions to the two equations
$p^2-Dq^2=\pm1$. In particular:
\begin{itemize}
 \item If $L$ is even, $+1$ occurs at every index $mL-1$, and the
 equation with right-hand side $-1$ has no integer solution.
 \item If $L$ is odd, $-1$ occurs for odd $m$ and $+1$ for even $m$.
\end{itemize}
The least positive solution of the $+1$ equation is
$(p_{L-1},q_{L-1})$ for even $L$ and $(p_{2L-1},q_{2L-1})$ for odd $L$.
\end{theorem}
\begin{proof}
All $Q_n$ with $n\ge1$ are positive. The lemma says that a convergent
solves one of the two equations exactly when $Q_n=1$.
By \eqref{eq:pell-state-return}, this happens exactly when $n=mL$;
\eqref{eq:pell-period} and its sign follow. Legendre's criterion, already
applied above, ensures that no positive solution lies outside the
convergents. Finally, the recurrence makes $p_0,p_1,\ldots$ strictly
increasing, so the first admissible index gives the least $p$, and hence
the least $q$ as well. If a solution of the $-1$ equation had either
coordinate zero, it would require $p^2=-1$ or $Dq^2=1$, both impossible;
changing signs reduces every other integer solution to a positive one.
\end{proof}

\begin{warningbox}[title={Stop before the final digit}]
The first candidate is
\[
 \frac{p_{L-1}}{q_{L-1}}=[a_0;a_1,\ldots,a_{L-1}],
\]
which stops \emph{before} the final $2a_0$. There are $L$ digits here,
counting $a_0$. For $\sqrt3=[1;\overline{1,2}]$, the candidate is
$[1;1]=2$, and $2^2-3\cdot1^2=1$. Including the next digit gives
$[1;1,2]=5/3$, whose remainder is $25-27=-2$.
\end{warningbox}

\begin{example}[An odd period: $D=13$]
The state calculation gives
$\sqrt{13}=[3;\overline{1,1,1,1,6}]$, with $L=5$.
The lemma explains every entry of the last column below.
\begin{center}
\begin{tabular}{ccrrr}
\toprule
$k$ & $a_k$ & $p_k/q_k$ & $Q_{k+1}$ & $p_k^2-13q_k^2$\\
\midrule
0 & 3 & $3/1$ & 4 & $-4$\\
1 & 1 & $4/1$ & 3 & $3$\\
2 & 1 & $7/2$ & 3 & $-3$\\
3 & 1 & $11/3$ & 4 & $4$\\
4 & 1 & $18/5$ & 1 & $-1$\\
5 & 6 & $119/33$ & 4 & $4$\\
6 & 1 & $137/38$ & 3 & $-3$\\
7 & 1 & $256/71$ & 3 & $3$\\
8 & 1 & $393/109$ & 4 & $-4$\\
9 & 1 & $649/180$ & 1 & $1$\\
\bottomrule
\end{tabular}
\end{center}
Thus the first negative solution is $(18,5)$, while the first positive
one is $(649,180)$. Squaring the former gives the latter:
\[
 (18+5\sqrt{13})^2=649+180\sqrt{13}.
\]
\end{example}

\subsection*{The period matrix becomes a Pell matrix}
Set $p=p_{L-1}$, $q=q_{L-1}$ and
$M=D(a_0)\cdots D(a_{L-1})$. At the returning state $P_L=a_0,Q_L=1$,
the two coefficient identities in the lemma give
\[
 M\begin{pmatrix}1&a_0\\0&1\end{pmatrix}
 =\begin{pmatrix}p&a_0p+p_{L-2}\\q&a_0q+q_{L-2}\end{pmatrix}
 =\begin{pmatrix}p&Dq\\q&p\end{pmatrix}.
\]
The shift by $a_0$ accounts for the fact that $M$ sends the tail
$\sqrt D+a_0$ to $\sqrt D$. The resulting matrix fixes $\sqrt D$ and
has determinant $p^2-Dq^2=(-1)^L$. For odd $L$, square it to obtain
determinant $+1$.

\subsection*{Why one fundamental solution generates all the others}
For $z=r+s\sqrt D$, define its \emph{norm} by
$N(z)=zz'=r^2-Ds^2$. Conjugation respects multiplication, so
$N(zw)=N(z)N(w)$. If $(u,v)$ is the least positive solution of the $+1$
equation, put $\varepsilon=u+v\sqrt D>1$. Its inverse is
$u-v\sqrt D$, again with integer coefficients.

\begin{proposition}[All positive Pell solutions are powers]
\label{prop:pell-powers}
Every positive solution of $p^2-Dq^2=1$ has the form
\[
 p+q\sqrt D=\varepsilon^j\qquad(j=1,2,\ldots).
\]
\end{proposition}
\begin{proof}
Powers have integer coefficients, norm $1$, and positive coefficients
for $j\ge1$. Conversely, let $\alpha=p+q\sqrt D>1$ have norm $1$.
Choose $j\ge0$ with $\varepsilon^j\le\alpha<\varepsilon^{j+1}$ and put
$\beta=\alpha\varepsilon^{-j}=r+s\sqrt D$. Then $r,s$ are integers,
$N(\beta)=1$, and $1\le\beta<\varepsilon$. Since $\beta'=1/\beta$,
\[
 r=\frac{\beta+\beta^{-1}}2>0,
 \qquad s=\frac{\beta-\beta^{-1}}{2\sqrt D}\ge0.
\]
If $\beta>1$, the function $t-t^{-1}$ is strictly increasing for $t>0$,
so $0<s<v$. This would be a smaller positive Pell solution, a
contradiction. Therefore $\beta=1$ and $\alpha=\varepsilon^j$.
Since $q>0$, we have $j\ge1$.
\end{proof}

Equivalently, multiplication by $\varepsilon$ is the matrix recurrence
\[
 \begin{pmatrix}p^{(j+1)}\\q^{(j+1)}\end{pmatrix}
 =\begin{pmatrix}u&Dv\\v&u\end{pmatrix}
  \begin{pmatrix}p^{(j)}\\q^{(j)}\end{pmatrix},
 \qquad (p^{(0)},q^{(0)})=(1,0).
\]
The superscript counts powers; the subscript on $p_k/q_k$ counts
continued-fraction convergents. For $D=2$, $\varepsilon=3+2\sqrt2$ gives
$(3,2),(17,12),(99,70)$.

\begin{bigidea}
Periodicity gives the exact return $Q_{mL}=1$. The determinant identity
then turns that return into $p_{mL-1}^2-Dq_{mL-1}^2=(-1)^{mL}$.
Legendre proves completeness, and powers of the least positive solution
give a second way to list every positive solution. Each step has a
different job.
\end{bigidea}

\begin{checkpoint}
Compute $99/70$ as a decimal. Its accuracy follows from
$99^2-2\cdot70^2=1$, which forces
\[
 \abs{\sqrt2-99/70}=\frac1{70^2(\sqrt2+99/70)}.
\]
For a second check, use the period of $\sqrt{23}$ above to find the
fundamental solution $(24,5)$ without searching through denominators.
\end{checkpoint}
